\scanpage{152}\mspage{7}%
A) Le formalisme des $Rf_!, f^!$ (combiné à $Rf_*, Lf^\alpha$, $\overset{\mathbb{L}}{\otimes}$ et $\mathbb{R}\mathrm{Hom}$, ``les six opérations'') porte implicitement en son sein la définition de l'homologie et l'identité essentielle entre \unsure{fonctions} homologiques et cohomologiques. \deleted{\unsure{E}} On dispose pour le moment de ce formalisme pour les faisceaux quasi-cohérents sur les schémas -- séminaire Hartshorne -- pour les espaces top. et les faisceaux de coeff. arbitraires dessus -- Verdier, exposé Bourbaki -- et en coh. étale des schémas pour les coefficients ``discrets'' \inserted{(de torsion ou ``$l$-adiques'')}
% Upper insertion guide: source (531,1024)--(1166,1140); immediately below the preceding addition, with its left descent beside the following lines. No arrowhead.
\figure{\begin{tikzpicture}[x=0.01cm,y=-0.01cm,line width=0.45pt,line cap=round,line join=round]
\draw (629,3) .. controls (633,11) and (632,13) .. (624,14) -- (41,18) .. controls (24,18) and (21,27) .. (21,42) -- (20,91) .. controls (20,104) and (17,112) .. (11,113);
\end{tikzpicture}}
\uline{premiers aux \deleted{\unsure{la|le}} car. résiduelles} -- SGA 5 --
% Right insertion guide and its interlinear semicolon: source (1545,1050)--(1645,1128). The punctuation stays inside the guide, above the end of the analytic addition; it is not duplicated in baseline prose.
\figure{\begin{tikzpicture}[x=0.025cm,y=-0.025cm,line width=0.45pt,line cap=round,line join=round]
\draw (95,6) -- (95,24) .. controls (94,30) and (86,33) .. (75,35) -- (30,38) .. controls (12,39) and (5,49) .. (3,61) -- (1,75);
\node[inner sep=0pt] at (29,59) {\inserted{;}};
\end{tikzpicture}}
\inserted{enfin pour les faisc. coh. sur les espaces analytiques -- Verdier-\unsure{Ruguet|Ruget} --}
(Le formalisme reste à développer dans le contexte cristallin, et en car. $0$ dans le contexte des modules stratifiés \inserted{à singularités}
% Guide beneath the preceding interlinear addition: source (1410,1322)--(1647,1379); the left curl descends before the baseline words ``à la''.
\figure{\begin{tikzpicture}[x=0.015cm,y=-0.015cm,line width=0.45pt,line cap=round,line join=round]
\draw (232,26) .. controls (180,27) and (116,9) .. (76,7) .. controls (52,6) and (26,12) .. (24,29) .. controls (23,39) and (24,49) .. (29,55);
\end{tikzpicture}}
à la Deligne, avec éventuellement -- sur le corps $\mathbb{C}$ -- des structures de Hodge à la clef, enfin dans le contexte des motifs ; je suis persuadé qu'il est présent
% An isolated short vertical spot below the initial u of the next ``un'' is recorded at source (243,1851)--(252,1870); it has no recoverable baseline punctuation position and is treated as a copying/stray mark.
un peu partout encore -- partout, peut-\unsure{\={e}tre|être}, où il y a un formalisme de nature cohomologique.) Travaillant en coh. étale sur un schéma \inserted{séparé}
% Guide beneath séparé: source (1063,2090)--(1175,2156); its left foot falls in the gap before ``de''.
\figure{\begin{tikzpicture}[x=0.025cm,y=-0.025cm,line width=0.45pt,line cap=round,line join=round]
\draw (99,6) .. controls (112,3) and (113,11) .. (101,18) .. controls (78,31) and (47,33) .. (24,40) .. controls (4,47) and (1,53) .. (8,64);
\end{tikzpicture}}
de t.f. sur un corps $k$, disons, avec un anneau de coeff. $\Lambda$ de torsion premier à \unsure{la} \unsure{car.|car}
% The image edge removes the bottom of the last words and any following punctuation; no conjectural comma is supplied.
% archive sheet number: 152
